\documentclass[11pt]{article}

\usepackage[T1]{fontenc}
\usepackage{lmodern}
\usepackage{amsmath,amssymb,amsthm}
\usepackage[margin=1in]{geometry}
\usepackage{enumitem}
\usepackage[colorlinks=true,linkcolor=blue,citecolor=blue,urlcolor=blue]{hyperref}

% ---------- LEAN NOTES: packages and macros used only in the section on the Lean formalization
\usepackage{upquote}
% \lean{name} typesets a Lean name; the line may break after a dot or an underscore.
\DeclareUrlCommand\lean{\urlstyle{tt}}
% The listings are copied from the Lean files and contain Unicode characters. Each of them is
% declared here; \LeanChar makes a symbol at least as wide as one typewriter character. The
% declarations are for pdfLaTeX (Overleaf's default compiler), not for XeLaTeX or LuaLaTeX.
\newsavebox{\LeanBox}
\newcommand{\LeanChar}[1]{\leavevmode\sbox{\LeanBox}{\ensuremath{#1}}%
  \ifdim\wd\LeanBox<\fontdimen2\font
    \makebox[\fontdimen2\font][c]{\usebox{\LeanBox}}%
  \else
    \usebox{\LeanBox}%
  \fi}
\DeclareUnicodeCharacter{03B9}{\LeanChar{\iota}}
\DeclareUnicodeCharacter{03BC}{\LeanChar{\mu}}
\DeclareUnicodeCharacter{211D}{\LeanChar{\mathbb{R}}}
\DeclareUnicodeCharacter{2102}{\LeanChar{\mathbb{C}}}
\DeclareUnicodeCharacter{2200}{\LeanChar{\forall}}
\DeclareUnicodeCharacter{2203}{\LeanChar{\exists}}
\DeclareUnicodeCharacter{2208}{\LeanChar{\in}}
\DeclareUnicodeCharacter{2209}{\LeanChar{\notin}}
\DeclareUnicodeCharacter{2227}{\LeanChar{\wedge}}
\DeclareUnicodeCharacter{2264}{\LeanChar{\le}}
\DeclareUnicodeCharacter{2260}{\LeanChar{\ne}}
\DeclareUnicodeCharacter{2192}{\LeanChar{\to}}
\DeclareUnicodeCharacter{2211}{\LeanChar{\Sigma}}
\DeclareUnicodeCharacter{2022}{\LeanChar{\bullet}}
\DeclareUnicodeCharacter{207B}{\LeanChar{{}^{-}}}
\DeclareUnicodeCharacter{00B9}{\LeanChar{{}^{1}}}
\DeclareUnicodeCharacter{1D65}{\LeanChar{{}_{v}}}
\DeclareUnicodeCharacter{2080}{\LeanChar{{}_{0}}}
\DeclareUnicodeCharacter{1D4DD}{\LeanChar{\mathcal{N}}}
% ---------- end of the LEAN NOTES block ----------

% ---------- theorem-like environments ----------
\newtheorem{theorem}{Theorem}
\newtheorem{proposition}[theorem]{Proposition}
\newtheorem{lemma}[theorem]{Lemma}
\newtheorem{corollary}[theorem]{Corollary}
\theoremstyle{definition}
\newtheorem{definition}[theorem]{Definition}
\newtheorem{example}[theorem]{Example}
\theoremstyle{remark}
\newtheorem{remark}[theorem]{Remark}

% ---------- macros ----------
\newcommand{\R}{\mathbb{R}}
\newcommand{\C}{\mathbb{C}}
\newcommand{\one}{\mathbf{1}}
\newcommand{\T}{\top}
\renewcommand{\Re}{\operatorname{Re}}
\newcommand{\diag}{\operatorname{diag}}
\newcommand{\cbar}{\bar{c}}

\setlist[enumerate,1]{label=\textup{(\alph*)}, leftmargin=2.2em}

\title{Convergence of the normalized state vector \\ in the positive linear system $\dot{x} = Mx$}
\author{Thomas M. English\footnote{Correspondence: Tom.English.PhD@gmail.com} \and Claude Fable 5.2 (Anthropic)}
\date{}

\begin{document}
\maketitle

\noindent
Let $M$ be a Metzler matrix and let $x(t)$ solve $\dot{x} = Mx$ from a nonnegative, nonzero
initial state. These notes prove that the normalized solution $x(t)/\sum_i x_i(t)$ converges to a
nonnegative eigenvector of $M$, with no irreducibility assumption
(Theorem~\ref{thm:main}).
% LEAN NOTES: the next sentence.
The proof has been formalized in the proof assistant Lean; Section~\ref{sec:lean} describes the
formalization.

% =====================================================================
\section{Setting and main result}\label{sec:setting}
% =====================================================================

\paragraph{Notation.}
For a real vector or matrix $A$ we write $A \ge 0$ if every entry of $A$ is nonnegative, and
$A \gg 0$ if every entry is positive. $I$ is the identity matrix, $\one$ is the column vector of
ones (of the appropriate dimension), and $e_j$ is the $j$th standard basis vector. For an index set $R \subseteq \{1,\dots,n\}$,
$x_R = (x_i)_{i \in R}$ is the corresponding subvector of $x$, and $M_{RR} = (M_{ij})_{i,j \in R}$
is the corresponding principal submatrix of $M$; in both, the indices of $R$ are kept in ascending
order. Eigenvalues may be complex, and $\Re\lambda$ is
the real part of $\lambda$.

\paragraph{Standing hypotheses.}
$M$ is a real $n \times n$ matrix with $M_{ij} \ge 0$ whenever $i \ne j$ (a \emph{Metzler} matrix).
The initial state $x(0) \in \R^n$ satisfies $x(0) \ge 0$ and $x(0) \ne 0$, and $x(t)$, $t \ge 0$,
is the solution of $\dot{x} = Mx$. We write
\[
  W(t) = \sum_{i=1}^n x_i(t) = \one^\T x(t),
  \qquad
  p(t) = \frac{x(t)}{W(t)} .
\]
If $x_i$ is the abundance of type $i$, then $W$ is the total population and $p$ is the vector of
type frequencies.

\begin{definition}[Reachability]\label{def:reach}
There is an \emph{arc} from $j$ to $i$ if $i \ne j$ and $M_{ij} > 0$. Index $i$ is
\emph{reachable from} $j$ if $i = j$, or if there are indices $j = i_0, i_1, \dots, i_r = i$ with
an arc from $i_{s-1}$ to $i_s$ for each $s = 1,\dots,r$. We set
\begin{align*}
  S &= \{\, j : x_j(0) > 0 \,\} && \text{(the support of the initial state),}\\
  R &= \{\, i : i \text{ is reachable from some } j \in S \,\} && \text{(the reachable set; } S \subseteq R\text{),}\\
  B &= M_{RR} && \text{(the matrix of the reduced system).}
\end{align*}
\end{definition}

An arc from $j$ to $i$ means that type $j$ contributes directly to the growth of type $i$.

\begin{theorem}\label{thm:main}
Under the standing hypotheses:
\begin{enumerate}
  \item\label{thm:main:support} For every $t > 0$, $x_i(t) > 0$ if $i \in R$; and $x_i(t) = 0$ for
        all $t \ge 0$ if $i \notin R$. In particular $W(t) > 0$ for all $t \ge 0$, so $p(t)$ is
        well defined.
  \item\label{thm:main:mu} The number
        $\mu = \max\{\, \Re\lambda : \lambda \text{ is an eigenvalue of } B \,\}$
        is itself an eigenvalue of $B$; in particular it is real.
  \item\label{thm:main:limit} The limit $p^* = \lim_{t \to \infty} p(t)$ exists, and
        \[
          p^* \ge 0, \qquad \one^\T p^* = 1, \qquad p^*_i = 0 \ \text{ for } i \notin R,
          \qquad M p^* = \mu\, p^* .
        \]
        In particular $\mu$ is an eigenvalue of $M$, with the nonnegative eigenvector $p^*$.
\end{enumerate}
\end{theorem}

The eigenvalue $\mu$ belongs to the reduced matrix $B$. It need not be the eigenvalue of $M$ with
the largest real part (Example~\ref{ex:reduction}).

% =====================================================================
\section{Background}\label{sec:background}
% =====================================================================

The proof uses three standard facts.

\begin{lemma}[The matrix exponential]\label{lem:exp}
For a real or complex square matrix $A$ and real $t$, the series
$e^{At} = \sum_{k=0}^{\infty} \frac{t^k}{k!} A^k$ converges, and:
\begin{enumerate}
  \item\label{lem:exp:ode} $\frac{d}{dt} e^{At} = A e^{At}$ and $e^{A \cdot 0} = I$. Hence the
        solution of $\dot{x} = Ax$ with initial value $x(0)$ is $x(t) = e^{At} x(0)$.
  \item\label{lem:exp:semigroup} $e^{A(s+t)} = e^{As} e^{At}$ for all real $s$, $t$.
  \item\label{lem:exp:shift} $e^{(A + cI)t} = e^{ct} e^{At}$ for every scalar $c \in \C$.
  \item\label{lem:exp:eig} If $Au = \lambda u$ with $\lambda \in \C$, then $e^{At} u = e^{\lambda t} u$.
\end{enumerate}
\end{lemma}

\begin{lemma}[Entries of the matrix exponential]\label{lem:modes}
Let $A$ be an $m \times m$ matrix with distinct eigenvalues $\lambda_1, \dots, \lambda_r \in \C$,
where $\lambda_\ell$ has algebraic multiplicity $m_\ell$. Then there are constant matrices
$C_{\ell,j} \in \C^{m \times m}$ such that, for all real $t$,
\begin{equation}\label{eq:modes}
  e^{At} = \sum_{\ell=1}^{r} \sum_{j=0}^{m_\ell - 1} t^{j} e^{\lambda_\ell t}\, C_{\ell,j} .
\end{equation}
In words: every entry of $e^{At}$ is a finite sum of terms $c\, t^{j} e^{\lambda t}$, with
$\lambda$ an eigenvalue of $A$ and $j$ smaller than the multiplicity of $\lambda$.
\end{lemma}

\begin{proof}
For $\lambda \in \C$ let
\[
  G_\lambda = \{\, v \in \C^{m} : (A - \lambda I)^{k} v = 0 \text{ for some } k \ge 0 \,\}
\]
be the \emph{generalized eigenspace} of $\lambda$. It is a subspace of $\C^{m}$ that $A$ maps into
itself, and it is not $\{0\}$ exactly when $\lambda$ is an eigenvalue of $A$. We use two standard
facts about generalized eigenspaces
% NOTE: Chapter 8 of this edition ("Operators on Complex Vector Spaces") treats generalized
% eigenspaces. (As far as I know, the book defines the multiplicity of an eigenvalue as the
% dimension of its generalized eigenspace and identifies the resulting characteristic polynomial
% with det(zI - A) in a later chapter.) The book itself was not inspected, hence "see, e.g.".
(see, e.g., \cite[Chapter~8]{Axler}):
\begin{enumerate}[label=\textup{(\roman*)}]
  \item every vector of $\C^{m}$ is a sum $v_1 + \dots + v_r$ with $v_\ell \in G_{\lambda_\ell}$
        for each $\ell$;
  \item the dimension of $G_{\lambda_\ell}$ is the algebraic multiplicity $m_\ell$, that is, the
        multiplicity of $\lambda_\ell$ as a root of the characteristic polynomial of $A$.
\end{enumerate}

\emph{The power $m_\ell$ of $A - \lambda_\ell I$ vanishes on $G_{\lambda_\ell}$.} The matrix
$N = A - \lambda_\ell I$ maps $G_{\lambda_\ell}$ into itself, and some power of $N$ vanishes on
$G_{\lambda_\ell}$: take a power that kills every vector of a basis. So the restriction of $N$ to
$G_{\lambda_\ell}$ is nilpotent. Its characteristic polynomial is therefore $z^{d}$, where $d$ is
the dimension of $G_{\lambda_\ell}$, and by the Cayley--Hamilton theorem $N^{d}$ vanishes on
$G_{\lambda_\ell}$. By (ii), $d = m_\ell$. Hence
\begin{equation}\tag{$\ast$}\label{eq:nilpotent}
  (A - \lambda_\ell I)^{m_\ell}\, v = 0 \qquad \text{for every } v \in G_{\lambda_\ell} .
\end{equation}

\emph{The exponential on $G_{\lambda_\ell}$.} Let $v \in G_{\lambda_\ell}$. By
Lemma~\ref{lem:exp}\ref{lem:exp:shift}, $e^{At} = e^{\lambda_\ell t}\, e^{(A - \lambda_\ell I)t}$.
Apply the series for $e^{(A - \lambda_\ell I)t}$ to $v$; by \eqref{eq:nilpotent} the terms with
$j \ge m_\ell$ vanish, so
\[
  e^{At} v = e^{\lambda_\ell t} \sum_{j=0}^{m_\ell - 1} \frac{t^{j}}{j!}\, (A - \lambda_\ell I)^{j} v .
\]

\emph{The columns of $e^{At}$.} By (i), write the $b$th standard basis vector as
$e_b = v_{b,1} + \dots + v_{b,r}$ with $v_{b,\ell} \in G_{\lambda_\ell}$. The $b$th column of
$e^{At}$ is
\[
  e^{At} e_b = \sum_{\ell=1}^{r} \sum_{j=0}^{m_\ell - 1} t^{j} e^{\lambda_\ell t}\,
  \frac{1}{j!}\, (A - \lambda_\ell I)^{j} v_{b,\ell} .
\]
This is \eqref{eq:modes}, with $C_{\ell,j}$ the matrix whose $b$th column is
$\frac{1}{j!} (A - \lambda_\ell I)^{j} v_{b,\ell}$, for $b = 1, \dots, m$.
\end{proof}

For example, $A = \begin{bmatrix} 1 & 1 \\ 0 & 1 \end{bmatrix}$ gives
$e^{At} = \begin{bmatrix} e^{t} & t e^{t} \\ 0 & e^{t} \end{bmatrix}$. Powers of $t$ occur only
when a repeated eigenvalue has too few independent eigenvectors; if $A$ is diagonalizable, only
$j = 0$ occurs.

\begin{lemma}[Perron--Frobenius for Metzler matrices]\label{lem:PF}
Let $B$ be a Metzler matrix and let $\mu$ be the largest real part of an eigenvalue of $B$. Then
$\mu$ is an eigenvalue of $B$, and every eigenvalue $\lambda \ne \mu$ of $B$ satisfies
$\Re\lambda < \mu$. The eigenvalue $\mu$ itself may be repeated, and no irreducibility is assumed.
\end{lemma}

\begin{proof}
For real vectors $u$ and $v$ we write $u \ge v$ if $u - v \ge 0$, and $u \le 0$ if $-u \ge 0$. For
a complex vector $z$, $|z|$ is the real vector with entries $|z_i|$.

\emph{\textup{(i)} If $A$ is a Metzler matrix, then $e^{At} \ge 0$ for $t \ge 0$.} Choose
$c \ge 0$ with $P = A + cI \ge 0$. By Lemma~\ref{lem:exp}\ref{lem:exp:shift},
$e^{At} = e^{-ct} e^{Pt} = e^{-ct} \sum_{k=0}^{\infty} \frac{t^{k}}{k!} P^{k}$, and every power
$P^{k}$ is a nonnegative matrix. (Step~1 of Section~\ref{sec:proof} repeats this argument and
refines it.)

\emph{\textup{(ii)} Inverse positivity.} Let $A$ be a Metzler matrix all of whose eigenvalues have
negative real part, and let $d$ be a real vector with $Ad \le 0$. Then $d \ge 0$.

To see this, put $\psi(t) = e^{At} d$. Since $A$ commutes with $e^{At}$,
Lemma~\ref{lem:exp}\ref{lem:exp:ode} gives $\dot{\psi}(t) = e^{At} A d$, and by (i) this vector is
$\le 0$ for $t \ge 0$. So every entry of $\psi$ is nonincreasing on $[0, \infty)$. By
Lemma~\ref{lem:modes}, every entry of $e^{At}$ is a sum of terms $a\, t^{j} e^{\lambda t}$ with
$\Re\lambda < 0$. Each such term tends to $0$ as $t \to \infty$, because its modulus is
$|a|\, t^{j} e^{(\Re\lambda) t}$ and exponential decay beats any power of $t$. Hence
$\psi(t) \to 0$. A nonincreasing function on $[0, \infty)$ that tends to $0$ is nonnegative, so
$d = \psi(0) \ge 0$.

\emph{\textup{(iii)} A comparison.} Let $\sigma$ be a real number with $\Re\lambda \le \sigma$ for
every eigenvalue $\lambda$ of $B$. Suppose that $Bw \ge r w$ for a real number $r$ and a real
vector $w \ge 0$ with $w \ne 0$. Then $r \le \sigma$.

Suppose instead that $r > \sigma$. The matrix $A = B - rI$ is Metzler, and its eigenvalues are the
numbers $\lambda - r$ with $\lambda$ an eigenvalue of $B$, so they have negative real part. Also
$A(-w) = -(Bw - rw) \le 0$. By (ii), $-w \ge 0$. With $w \ge 0$ this gives $w = 0$, a
contradiction.

\emph{\textup{(iv)} Moduli of an eigenvector.} Let $Bz = \lambda z$ with $\lambda \in \C$ and $z$
a complex vector, and choose $c \ge 0$ with $P = B + cI \ge 0$. Then $Pz = (\lambda + c) z$, so for
every index $i$
\[
  |\lambda + c|\, |z_i| = \Bigl| \sum_{j} P_{ij} z_j \Bigr| \le \sum_{j} P_{ij}\, |z_j|
  = (B|z|)_i + c\, |z_i| .
\]
That is, $B|z| \ge (|\lambda + c| - c)\, |z|$.

\emph{\textup{(v)} Conclusion.} Let $\lambda$ be an eigenvalue of $B$ with $\Re\lambda = \mu$, and
let $z \ne 0$ be an eigenvector for it. The vector $|z| \ge 0$ is not zero, so (iv) and (iii), with
$\sigma = \mu$, $w = |z|$ and $r = |\lambda + c| - c$, give $|\lambda + c| - c \le \mu$. Hence
\[
  |\lambda + c| \le \mu + c = \Re(\lambda + c) .
\]
A complex number whose modulus is at most its real part is real. So $\lambda$ is real, that is,
$\lambda = \mu$. By the definition of $\mu$, some eigenvalue of $B$ has real part $\mu$; we have
just seen that it equals $\mu$. So $\mu$ is an eigenvalue of $B$, and it is the only eigenvalue
with real part $\mu$: every other eigenvalue has real part less than $\mu$.
\end{proof}

Lemma~\ref{lem:PF} is the form, for Metzler matrices, of the classical theorem that the spectral
radius of a nonnegative matrix is an eigenvalue of that matrix
% NOTE: section placement confirmed from the publisher's table of contents, and the number agrees
% with how other papers cite the 2013 edition; the page itself was not inspected.
\cite[Theorem~8.3.1]{HJ}. It also follows from that theorem, applied to $P = B + cI$.

% =====================================================================
\section{Proof of Theorem~\ref{thm:main}}\label{sec:proof}
% =====================================================================

\subsection*{Step 1. $e^{Mt}$ is nonnegative, and reachability decides which entries are positive}

Choose $c \ge 0$ with $P = M + cI \ge 0$; for instance $c = \max_i |M_{ii}|$. By
Lemma~\ref{lem:exp}\ref{lem:exp:shift},
\[
  e^{Mt} = e^{-ct} e^{Pt} = e^{-ct} \sum_{k=0}^{\infty} \frac{t^{k}}{k!} P^{k} .
\]
Every power $P^k$ is a nonnegative matrix, so $e^{Mt} \ge 0$ for $t \ge 0$. Moreover, for $t > 0$
the $(i,j)$ entry of $e^{Mt}$ is positive if $[P^k]_{ij} > 0$ for at least one $k \ge 0$, and it is
zero for all $t \ge 0$ if $[P^k]_{ij} = 0$ for every $k \ge 0$.

For $k \ge 1$, the entry $[P^k]_{ij}$ is the sum, over all sequences of indices $j = i_0, i_1, \dots, i_k = i$,
of the products
\[
  P_{i_k i_{k-1}} \cdots P_{i_2 i_1} P_{i_1 i_0} .
\]
These products are nonnegative, so the sum is positive exactly when some sequence has every factor
positive. In such a sequence each step
either stays put ($i_s = i_{s-1}$, a diagonal entry of $P$) or follows an arc
($i_s \ne i_{s-1}$ and $M_{i_s i_{s-1}} > 0$, since $P$ and $M$ have the same off-diagonal
entries). Deleting the steps that stay put shows that $i$ is reachable from $j$. Conversely, if $i$
is reachable from $j$ along $k \ge 1$ arcs then $[P^k]_{ij} > 0$, and the case $i = j$ is covered by
$P^0 = I$. Therefore, for every $t > 0$,
\begin{equation}\label{eq:pattern}
  [e^{Mt}]_{ij} > 0 \ \text{ if $i$ is reachable from $j$,}
  \qquad
  [e^{Mt}]_{ij} = 0 \ \text{ otherwise.}
\end{equation}

\subsection*{Step 2. Only the reachable indices matter}

Since $x(t) = e^{Mt} x(0)$ and $x_j(0) = 0$ for $j \notin S$,
\[
  x_i(t) = \sum_{j \in S} [e^{Mt}]_{ij}\, x_j(0) ,
\]
a sum of nonnegative terms. If $i \in R$, some $j \in S$ has $i$ reachable from $j$, and by
\eqref{eq:pattern} the corresponding term is positive for $t > 0$. If $i \notin R$, every term is
zero. Hence, for every $t > 0$,
\begin{equation}\label{eq:support}
  x_i(t) > 0 \ \text{ if } i \in R,
  \qquad
  x_i(t) = 0 \ \text{ if } i \notin R,
\end{equation}
and $x_i(t) = 0$ for all $t \ge 0$ when $i \notin R$. Three consequences follow.

\begin{itemize}[leftmargin=2em]
  \item \emph{The total is positive.} $W(0) > 0$ because $x(0) \ge 0$ is nonzero, and $W(t) > 0$
        for $t > 0$ by \eqref{eq:support}. This proves
        part~\ref{thm:main:support} of the theorem.
  \item \emph{No arc leaves $R$.} If $j \in R$ and $i \notin R$, then
        \begin{equation}\label{eq:noarc}
          M_{ij} = 0 ,
        \end{equation}
        because otherwise $i$ would be reachable from $j$, and hence from $S$.
  \item \emph{The reduced system.} Put $y(t) = x_R(t)$. For $i \in R$,
        $\dot{x}_i = \sum_{j} M_{ij} x_j = \sum_{j \in R} M_{ij} x_j$, because $x_j \equiv 0$ for
        $j \notin R$. Hence
        \[
          \dot{y} = By, \qquad y(t) = e^{Bt} y(0), \qquad y(t) \gg 0 \ \text{ for every } t > 0 .
        \]
        The matrix $B$ is Metzler, because its off-diagonal entries are off-diagonal entries of $M$.
\end{itemize}

\subsection*{Step 3. The leading term of $e^{Bt}$}

Apply Lemma~\ref{lem:PF} to $B$. This proves part~\ref{thm:main:mu} of the theorem, and shows that
every eigenvalue $\lambda \ne \mu$ of $B$ has $\Re\lambda < \mu$. Apply Lemma~\ref{lem:modes} to
$B$ and write $C_j$ for the coefficient matrix of $t^{j} e^{\mu t}$:
\begin{equation}\label{eq:expansion}
  e^{Bt} = \sum_{j} t^{j} e^{\mu t}\, C_j
  \;+\; \sum_{\lambda \ne \mu} \sum_{j} t^{j} e^{\lambda t}\, C_{\lambda,j} .
\end{equation}

\emph{Some $C_j$ is nonzero.} Take $u \ne 0$ with $Bu = \mu u$. By
Lemma~\ref{lem:exp}\ref{lem:exp:eig}, $e^{-\mu t} e^{Bt} u = u$ for all $t$. If every $C_j$ were
zero, then by \eqref{eq:expansion} each entry of $e^{-\mu t} e^{Bt}$ would be a sum of terms
$c\, t^{j} e^{(\lambda - \mu)t}$ with $\Re\lambda < \mu$. Each such term tends to $0$ as
$t \to \infty$, because its modulus is $|c|\, t^{j} e^{-(\mu - \Re\lambda)t}$ and exponential decay
beats any power of $t$. This would give $u = 0$, a contradiction.

\emph{Definition of $k$ and $L$.} Let $k$ be the largest $j$ with $C_j \ne 0$, and put $L = C_k$.
Divide \eqref{eq:expansion} by $t^{k} e^{\mu t}$. Every term other than $L$ carries a factor
$t^{j-k}$ with $j < k$, or a factor $t^{j-k} e^{(\lambda - \mu)t}$ with $\Re\lambda < \mu$; both
tend to $0$. Therefore
\begin{equation}\label{eq:L}
  \lim_{t \to \infty} \frac{e^{Bt}}{t^{k} e^{\mu t}} = L \ne 0 .
\end{equation}
The fastest-growing terms of $e^{Bt}$ grow like $t^{k} e^{\mu t}$, and $L$ holds their
coefficients. Both $k$ and $L$ depend only on $B$.

\subsection*{Step 4. $L$ is nonnegative, and its nonzero columns are eigenvectors}

\emph{\textup{(i)} $L \ge 0$.} By Step~1 applied to the Metzler matrix $B$, $e^{Bt} \ge 0$; and
$t^{k} e^{\mu t} > 0$ for $t > 0$. A limit of nonnegative matrices is nonnegative; in particular
$L$ is real.

\emph{\textup{(ii)} $BL = \mu L$.} Fix $s > 0$. Since $e^{Bs} e^{Bt} = e^{B(s+t)}$ by
Lemma~\ref{lem:exp}\ref{lem:exp:semigroup},
\[
  e^{Bs} \cdot \frac{e^{Bt}}{t^{k} e^{\mu t}}
  \;=\;
  \frac{(s+t)^{k} e^{\mu(s+t)}}{t^{k} e^{\mu t}}
  \cdot
  \frac{e^{B(s+t)}}{(s+t)^{k} e^{\mu(s+t)}} .
\]
Let $t \to \infty$. The left side tends to $e^{Bs} L$ by \eqref{eq:L}. On the right, the first
factor equals $(1 + s/t)^{k} e^{\mu s}$ and tends to $e^{\mu s}$, and the second factor tends to
$L$ by \eqref{eq:L}. Hence
\begin{equation}\label{eq:semigroup}
  e^{Bs} L = e^{\mu s} L \qquad \text{for every } s > 0 .
\end{equation}
Subtract $L$ from both sides and divide by $s$:
\[
  \frac{e^{Bs} - I}{s}\, L = \frac{e^{\mu s} - 1}{s}\, L .
\]
As $s \to 0^{+}$ the left side tends to $BL$, because
$\frac{d}{ds} e^{Bs} \big|_{s=0} = B$ by Lemma~\ref{lem:exp}\ref{lem:exp:ode}, and the right side
tends to $\mu L$. Therefore
\begin{equation}\label{eq:eig}
  BL = \mu L .
\end{equation}
Consequently, for every vector $z$ we have $B(Lz) = \mu\,(Lz)$: the vector $Lz$ is either zero or
an eigenvector of $B$ for the eigenvalue $\mu$.

\emph{\textup{(iii)} If $z \gg 0$, then $Lz \ne 0$.} By \eqref{eq:L} and (i), some entry $L_{ij}$
is positive. Then $(Lz)_i = \sum_{l} L_{il} z_l \ge L_{ij} z_j > 0$.

\subsection*{Step 5. The limit for the reduced system}

Fix any $\tau > 0$ and put $z = y(\tau)$; by Step~2, $z \gg 0$. For $t > 0$,
Lemma~\ref{lem:exp}\ref{lem:exp:semigroup} gives $y(\tau + t) = e^{Bt} z$, so by \eqref{eq:L}
\[
  \frac{y(\tau + t)}{t^{k} e^{\mu t}} \;\longrightarrow\; v := Lz
  \qquad (t \to \infty).
\]
By Step~4, $v \ge 0$, $v \ne 0$ and $Bv = \mu v$. Summing the entries,
$\one^\T y(\tau + t) / (t^{k} e^{\mu t}) \to \one^\T v > 0$. Dividing the two limits, which is
legitimate because the limit of the denominator is not zero,
\begin{equation}\label{eq:q}
  \frac{y(t)}{\one^\T y(t)} \;\longrightarrow\; q := \frac{v}{\one^\T v}
  \qquad (t \to \infty),
  \qquad\text{with}\quad q \ge 0, \quad \one^\T q = 1, \quad Bq = \mu q .
\end{equation}

\subsection*{Step 6. Back to the full system}

By \eqref{eq:support}, $W(t) = \one^\T y(t)$. Hence $p_R(t) = y(t)/\one^\T y(t)$, and $p_i(t) = 0$
for $i \notin R$. By \eqref{eq:q}, $p(t) \to p^*$, where $p^*_R = q$ and $p^*_i = 0$ for
$i \notin R$. So $p^* \ge 0$ and $\one^\T p^* = 1$.

It remains to check $M p^* = \mu p^*$. For any $i$, $(M p^*)_i = \sum_{j \in R} M_{ij} q_j$.
If $i \in R$, this is $(Bq)_i = \mu q_i = \mu p^*_i$. If $i \notin R$, every $M_{ij}$ in the sum is
zero by \eqref{eq:noarc}, so $(M p^*)_i = 0 = \mu p^*_i$. This proves
part~\ref{thm:main:limit}. \qed

\begin{remark}[What each hypothesis does]\label{rem:roles}
The Metzler property is used in four places. It gives $e^{Mt} \ge 0$ and, because sums of
nonnegative products cannot cancel, the exact positivity pattern \eqref{eq:pattern} (Step~1). It
makes the absence of an arc from $j$ to $i$ mean $M_{ij} = 0$, which is \eqref{eq:noarc} (Step~2).
It makes the eigenvalue of largest real part real and strictly dominant, so that the leading term
of $e^{Bt}$ does not oscillate (Lemma~\ref{lem:PF}, Step~3). And it gives $L \ge 0$ (Step~4). Strict positivity of
$y(\tau)$ is used once, in Step~5, to guarantee $L y(\tau) \ne 0$. In a general linear system both
can fail: nonreal eigenvalues cause persistent oscillation in the direction of the solution when
they contribute to its leading term, and an initial state can have no component along the leading
term of the matrix exponential. The reduction to $R$ is what supplies strict positivity. Applied
to the full matrix $M$, the argument would fail: in Example~\ref{ex:reduction} the leading
coefficient matrix of $e^{Mt}$ annihilates $x(0)$, and the solution grows at a rate smaller than
the largest eigenvalue of $M$.
\end{remark}

\begin{example}[The reduction to $R$ is needed]\label{ex:reduction}
Let $M = \begin{bmatrix} 1 & 0 \\ 0 & 2 \end{bmatrix}$. For $x(0) = (1, 0)^\T$ we get $R = \{1\}$,
$B = [\,1\,]$, $\mu = 1$ and $p(t) = e_1$ for all $t$, although the dominant eigenvalue of $M$ is
$2$. For $x(0) \gg 0$ we get $R = \{1,2\}$, $\mu = 2$ and $p^* = e_2$.
\end{example}

% >>>>>>>>>>>>>>>>>>>>>>>>>>>>>>>>>>>>>>>>>>>>>>>>>>>>>>>>>>>>>>>>>>>>>>>>>>>>>>>>>>>>>>>>>>>>>>>>>>>
% LEAN NOTES. This section, the block marked "LEAN NOTES" in the preamble and the sentence marked
% "LEAN NOTES" in the opening paragraph are the only places where this file mentions Lean. A
% version of the file without them is obtained by deleting those three pieces.
% >>>>>>>>>>>>>>>>>>>>>>>>>>>>>>>>>>>>>>>>>>>>>>>>>>>>>>>>>>>>>>>>>>>>>>>>>>>>>>>>>>>>>>>>>>>>>>>>>>>
% =====================================================================
\section{Notes on the Lean formalization}\label{sec:lean}
% =====================================================================

Theorem~\ref{thm:main} and the whole of its proof, that is, Lemmas~\ref{lem:exp}--\ref{lem:PF} and
Steps~1--6, have been formalized in the proof assistant Lean~4, on top of the mathematical library
Mathlib, and Lean accepts the formal proof. Remark~\ref{rem:roles}, Example~\ref{ex:reduction} and
the unnumbered comments between the numbered statements are commentary; they are not formalized.

This section says where the formal statements are, how to read the formal statement of the
theorem, and where the formal proof differs from the text. The Lean files are the authority: the
listings below are copied from them, and if a listing and a file ever disagree, the file is right.

\subsection{The files}\label{sec:lean:files}

The Lean project is called \lean{pos_linear_sys_convergence}. It uses Lean~4.36.0 and
Mathlib~v4.36.0. Its library \lean{PosLinearSysConvergence} consists of six files:
\begin{center}
\begin{tabular}{@{}ll@{}}
\hline
File & Content \\
\hline
\lean{Defs.lean} & Metzler matrices; Definition~\ref{def:reach} \\
\lean{MatrixExp.lean} & Lemma~\ref{lem:exp} \\
\lean{Positivity.lean} & Step~1 \\
\lean{Modes.lean} & Lemma~\ref{lem:modes}; Step~3 \\
\lean{PerronFrobenius.lean} & Lemma~\ref{lem:PF} \\
\lean{Main.lean} & Steps~2, 4, 5 and~6; Theorem~\ref{thm:main} \\
\hline
\end{tabular}
\end{center}
Every name defined in these files begins with the prefix \lean{PosLinearSys.}, which is omitted
below. So the full name of the main theorem is \lean{PosLinearSys.tendsto_normalized_solution}.

\subsection{The definitions}\label{sec:lean:defs}

The standing hypothesis on $M$ and Definition~\ref{def:reach} are the following five definitions,
from \lean{Defs.lean} (the explanatory comments in the file are left out).
{\small
\begin{verbatim}
def IsMetzler (M : Matrix ι ι ℝ) : Prop :=
  ∀ i j, i ≠ j → 0 ≤ M i j

def Arc (M : Matrix ι ι ℝ) (j i : ι) : Prop :=
  i ≠ j ∧ 0 < M i j

def Reachable (M : Matrix ι ι ℝ) (j i : ι) : Prop :=
  Relation.ReflTransGen (Arc M) j i

open Classical in
noncomputable def reachSet (M : Matrix ι ι ℝ) (x₀ : ι → ℝ) : Finset ι :=
  Finset.univ.filter fun i => ∃ j, 0 < x₀ j ∧ Reachable M j i

def reducedMatrix (M : Matrix ι ι ℝ) (x₀ : ι → ℝ) :
    Matrix (reachSet M x₀) (reachSet M x₀) ℝ :=
  M.submatrix (fun i => (i : ι)) (fun i => (i : ι))
\end{verbatim}
}
They are read as follows.

\begin{itemize}[leftmargin=2em]
  \item \emph{Indices, vectors and matrices.} The index set $\{1, \dots, n\}$ of the text is
        replaced by an arbitrary finite set \verb|ι| (in Lean, a finite \emph{type}). A vector is
        a function \verb|ι → ℝ|, so \verb|x₀ j| is the entry $(x_0)_j$. \verb|Matrix ι ι ℝ| is the
        type of real matrices with rows and columns indexed by \verb|ι|, and \verb|M i j| is the
        entry $M_{ij}$. The first three definitions make sense for any \verb|ι|. The last two need
        \verb|ι| to be finite, an assumption that the file declares on a separate line.
  \item \emph{Metzler matrices.} \verb|IsMetzler M| says that $M_{ij} \ge 0$ whenever $i \ne j$.
  \item \emph{Arcs.} \verb|Arc M j i| says that there is an arc from $j$ to $i$. The arguments are
        in the order source, target, whereas the matrix entry is \verb|M i j|, with the target as
        the row index, exactly as in Definition~\ref{def:reach}.
  \item \emph{Reachability.} \lean{Relation.ReflTransGen} is Mathlib's reflexive and transitive
        closure of a relation: \verb|Relation.ReflTransGen r a b| holds exactly when $b = a$ or
        there is a finite chain $a = c_0, c_1, \dots, c_r = b$ in which $c_{s-1}$ is related to
        $c_s$ by \verb|r| for each $s$. So \verb|Reachable M j i| says that $i$ is reachable
        from $j$ in the sense of Definition~\ref{def:reach}.
  \item \emph{The reachable set.} \verb|reachSet M x₀| is the set $R$ that belongs to the initial
        state $x_0$: the expression \lean{Finset.univ.filter} selects, from the set of all
        indices, those $i$ for which some $j$ has $(x_0)_j > 0$ and $i$ reachable from $j$. The
        support $S$ is not given a name. The words \verb|open Classical in| and
        \lean{noncomputable} record that membership in $R$ is settled by classical logic, not by
        an algorithm. They have no mathematical content.
  \item \emph{The reduced matrix.} \verb|reducedMatrix M x₀| is $B = M_{RR}$. Its rows and columns
        are indexed by the elements of $R$ themselves: \verb|M.submatrix f g| is the matrix with
        entries $M_{f(i)\,g(j)}$, and here $f$ and $g$ are both the inclusion of $R$ in the index
        set. No ordering of $R$ is involved, so the convention of Section~\ref{sec:setting} that
        the indices of $R$ are kept in ascending order has no counterpart, and none is needed.
\end{itemize}

\subsection{The main theorem}\label{sec:lean:main}

Theorem~\ref{thm:main} is the following statement, from \lean{Main.lean}. In the file it is
followed by its proof.
{\footnotesize
\begin{verbatim}
theorem tendsto_normalized_solution (M : Matrix ι ι ℝ) (hM : IsMetzler M) (x : ℝ → ι → ℝ)
    (hx0 : ∀ i, 0 ≤ x 0 i) (hx0' : x 0 ≠ 0)
    (hx : ∀ t ∈ Set.Ici (0 : ℝ), HasDerivWithinAt x (M *ᵥ x t) (Set.Ici 0) t) :
    ((∀ t : ℝ, 0 < t → ∀ i ∈ reachSet M (x 0), 0 < x t i) ∧
      (∀ t : ℝ, 0 ≤ t → ∀ i ∉ reachSet M (x 0), x t i = 0) ∧
      (∀ t : ℝ, 0 ≤ t → 0 < ∑ i, x t i)) ∧
    ∃ μ : ℝ,
      (IsGreatest
          (Complex.re '' spectrum ℂ ((reducedMatrix M (x 0)).map (algebraMap ℝ ℂ))) μ ∧
        μ ∈ spectrum ℝ (reducedMatrix M (x 0))) ∧
      ∃ p : ι → ℝ,
        Tendsto (fun t : ℝ => (∑ i, x t i)⁻¹ • x t) atTop (nhds p) ∧
          (∀ i, 0 ≤ p i) ∧ ∑ i, p i = 1 ∧ (∀ i ∉ reachSet M (x 0), p i = 0) ∧
            M *ᵥ p = μ • p ∧ μ ∈ spectrum ℝ M
\end{verbatim}
}
The statement is preceded in the file by the line
\begin{center}
\verb|variable {ι : Type*} [Fintype ι] [DecidableEq ι]|
\end{center}
which says that \verb|ι| is a finite type in which equality of two indices can be decided. The
second assumption is a technicality, needed to write down the identity matrix; in classical logic
it holds for every type.

\paragraph{Hypotheses.}
\begin{itemize}[leftmargin=2em]
  \item \verb|(M : Matrix ι ι ℝ) (hM : IsMetzler M)|: $M$ is a Metzler matrix. Each hypothesis has
        a name (here \lean{hM}), by which the proof refers to it.
  \item \verb|(x : ℝ → ι → ℝ)|: the solution, with \verb|x t i| standing for $x_i(t)$. The
        function is defined for every real $t$, but the hypotheses and the conclusions involve
        only $t \ge 0$. The values of $x$ at negative times are arbitrary and play no role.
  \item \verb|hx0| and \verb|hx0'|: $x(0) \ge 0$ and $x(0) \ne 0$.
  \item \verb|hx|: $x$ solves $\dot{x} = Mx$ on $[0, \infty)$. Here \verb|Set.Ici 0| is the
        interval $[0, \infty)$, \verb|M *ᵥ x t| is the product $M x(t)$ of a matrix and a vector,
        and \verb|HasDerivWithinAt x v (Set.Ici 0) t| says that the restriction of $x$ to
        $[0, \infty)$ has derivative $v$ at $t$. For $t > 0$ this is the ordinary derivative. At
        $t = 0$ it is the derivative from the right.
\end{itemize}

\paragraph{Conclusions.}
\begin{itemize}[leftmargin=2em]
  \item \emph{Part~\ref{thm:main:support}} is the first three lines after the hypotheses:
        $x_i(t) > 0$ for $t > 0$ and $i \in R$; $x_i(t) = 0$ for $t \ge 0$ and $i \notin R$; and
        $W(t) > 0$ for $t \ge 0$. The third clause, which the text presents as a consequence of
        the first two, is stated separately.
  \item \verb|∃ μ : ℝ| introduces the number $\mu$. Everything that follows is about this $\mu$,
        so parts~\ref{thm:main:mu} and~\ref{thm:main:limit} concern the same number.
  \item \emph{Part~\ref{thm:main:mu}} is the conjunction that begins with \lean{IsGreatest}.
        In it,
        \begin{center}
        \verb|(reducedMatrix M (x 0)).map (algebraMap ℝ ℂ)|
        \end{center}
        is $B$ regarded as a complex matrix. For a square matrix $A$, \verb|spectrum ℂ A| is the
        set of complex
        numbers $\lambda$ for which $\lambda I - A$ is not invertible, that is, the set of
        eigenvalues of $A$. \verb|Complex.re '' S| is the image of the set $S$ under $\Re$, and
        \verb|IsGreatest S μ| says that $\mu$ belongs to $S$ and that no element of $S$ exceeds
        it. So the first clause says that
        $\mu = \max\{\, \Re\lambda : \lambda \text{ is an eigenvalue of } B \,\}$, and the second,
        \verb|μ ∈ spectrum ℝ (reducedMatrix M (x 0))|, says that $\mu$ is an eigenvalue of the
        real matrix $B$. A set has at most one greatest element, so the $\mu$ whose existence is
        asserted is the number $\mu$ of the theorem.
  \item \emph{Part~\ref{thm:main:limit}.} \verb|∃ p : ι → ℝ| introduces the vector $p^*$.
        \verb|Tendsto f atTop (nhds p)| says that $f(t) \to p$ as $t \to \infty$, and
        \verb|(∑ i, x t i)⁻¹ • x t| is $p(t) = x(t)/W(t)$, so the first clause says that
        $p(t) \to p^*$. (Vectors converge when each entry does.) The remaining clauses are, in
        order: $p^* \ge 0$; $\one^\T p^* = 1$; $p^*_i = 0$ for $i \notin R$; $M p^* = \mu p^*$;
        and $\mu$ is an eigenvalue of $M$. The last clause is the ``in particular'' of the text,
        stated separately.
\end{itemize}

\paragraph{Three conventions.}
\begin{itemize}[leftmargin=2em]
  \item \emph{Division by zero.} In Lean the inverse of $0$ is defined, and equals $0$. So the
        formal expression for $p(t)$ makes sense for every $t$, whether or not $W(t) \ne 0$. By
        part~\ref{thm:main:support}, $W(t) > 0$ for $t \ge 0$, so for these $t$ it is the true
        quotient $x(t)/W(t)$; and the limit as $t \to \infty$ involves no other $t$.
  \item \emph{Eigenvector.} The statement does not use the word. It gives $M p^* = \mu p^*$ and
        $\one^\T p^* = 1$; the second makes $p^* \ne 0$, so $p^*$ is an eigenvector.
  \item \emph{The size $n$.} The statement does not mention $n$. If the index set is empty,
        there is no nonzero $x(0)$, so the hypotheses cannot be met; the theorem is about
        $n \ge 1$, as in the text.
\end{itemize}

\subsection{Three variants of the main theorem}\label{sec:lean:variants}

\lean{Main.lean} contains three further theorems. They are not in the text. Each has the
conclusions of the main theorem and differs only in how it says that $x$ is the solution of
$\dot{x} = Mx$. Each is derived from the main theorem.

\begin{itemize}[leftmargin=2em]
  \item \lean{tendsto_normalized_solution_of_hasDerivAt} replaces \verb|hx| by
        \begin{center}
        \verb|(hx : ∀ t : ℝ, 0 ≤ t → HasDerivAt x (M *ᵥ x t) t)|
        \end{center}
        that is, $x$ has the ordinary, two-sided derivative $M x(t)$ at every $t \ge 0$, including
        $t = 0$. This asks more of $x$ than the main theorem does.
  \item \lean{tendsto_normalized_solution_of_continuousOn} replaces \verb|hx| by
        \begin{center}
        \verb|(hc : ContinuousOn x (Set.Ici 0))| \\
        \verb|(hx : ∀ t : ℝ, 0 < t → HasDerivAt x (M *ᵥ x t) t)|
        \end{center}
        that is, $x$ is continuous on $[0, \infty)$ and has derivative $M x(t)$ at every $t > 0$.
        Nothing is assumed about a derivative at $t = 0$; the proof shows that such an $x$ has
        the derivative from the right $M x(0)$ there. Conversely, a function that satisfies
        \verb|hx| of the main theorem is continuous on $[0, \infty)$ and has the ordinary
        derivative $M x(t)$ at every $t > 0$. So the two hypotheses are equivalent.
  \item \lean{tendsto_normalized_solution_exp} has no differential equation among its hypotheses.
        It takes a vector $x_0 \ge 0$, $x_0 \ne 0$, and states the conclusions for the function
        $x(t) = e^{Mt} x_0$, written \verb|exp (t • M) *ᵥ x₀|.
\end{itemize}

The variants serve two purposes. First, ``$x$ is the solution on $[0, \infty)$'' can be made
precise in more than one reasonable way, and the first two variants show that the result does not
depend on the choice made in the main theorem. Second, a theorem whose hypotheses can never be
met is true and says nothing. The third variant shows that this is not the case here: for every
Metzler matrix $M$ and every nonnegative, nonzero $x_0$ there is a function $x$ that satisfies
the hypotheses of the main theorem and has $x(0) = x_0$.

\subsection{Where each part of the text is}\label{sec:lean:map}

The table lists, for each part of the text, the Lean theorem or definition that formalizes it.
The file is named in the heading of each group; a name from another file has its file in
parentheses.

% The table is set as one small table per group, so that a page break can fall between groups.
\par\smallskip
{\small
\noindent\begin{tabular}{@{}p{0.36\textwidth}p{0.60\textwidth}@{}}
\hline
Text & Lean \\
\hline
\multicolumn{2}{@{}l}{\emph{Section~\ref{sec:setting}} (\lean{Defs.lean})} \\
Metzler matrix & \lean{IsMetzler} \\
Definition~\ref{def:reach} & \lean{Arc}, \lean{Reachable}, \lean{reachSet}, \lean{reducedMatrix} \\
\end{tabular}\par
\vspace{0.4ex}
\noindent\begin{tabular}{@{}p{0.36\textwidth}p{0.60\textwidth}@{}}
\multicolumn{2}{@{}l}{\emph{Lemma~\ref{lem:exp}} (\lean{MatrixExp.lean})} \\
part~\ref{lem:exp:ode} & \lean{hasDerivAt_exp_smul}, \lean{exp_zero_smul} \\
part~\ref{lem:exp:ode}, ``Hence \dots'' & \lean{eq_exp_smul_mulVec_of_hasDerivWithinAt} \\
part~\ref{lem:exp:semigroup} & \lean{exp_add_smul} \\
part~\ref{lem:exp:shift} & \lean{exp_smul_add_smul_one}, \lean{exp_smul_eq_cexp_smul} \\
part~\ref{lem:exp:eig} & \lean{exp_smul_mulVec_of_mulVec_eq_smul} \\
\end{tabular}\par
\vspace{0.4ex}
\noindent\begin{tabular}{@{}p{0.36\textwidth}p{0.60\textwidth}@{}}
\multicolumn{2}{@{}l}{\emph{Lemma~\ref{lem:modes}} (\lean{Modes.lean})} \\
the lemma & \lean{exp_smul_eq_sum_modes} \\
$v \in G_\lambda$ & \verb|MemGenEigenspace A lam v| \\
fact (i) & \lean{exists_genEigenspace_decomp} \\
equation \eqref{eq:nilpotent}, with fact (ii) & \lean{MemGenEigenspace.pow_rootMultiplicity_mulVec} \\
the exponential on $G_{\lambda_\ell}$ & \lean{exp_smul_mulVec_of_pow_mulVec_eq_zero} (\lean{MatrixExp.lean}) \\
\end{tabular}\par
\vspace{0.4ex}
\noindent\begin{tabular}{@{}p{0.36\textwidth}p{0.60\textwidth}@{}}
\multicolumn{2}{@{}l}{\emph{Lemma~\ref{lem:PF}} (\lean{PerronFrobenius.lean})} \\
the lemma & \lean{exists_dominant_eigenvalue}, \lean{isGreatest_re_spectrum} \\
part (i) & \lean{exp_smul_apply_nonneg} (\lean{Positivity.lean}) \\
part (ii) & \lean{nonneg_of_mulVec_nonpos} \\
part (iii) & \lean{le_of_mul_le_mulVec} \\
part (iv) & \lean{norm_add_sub_mul_le_mulVec} \\
part (v) & the proof of \lean{exists_dominant_eigenvalue} \\
\end{tabular}\par
\vspace{0.4ex}
\noindent\begin{tabular}{@{}p{0.36\textwidth}p{0.60\textwidth}@{}}
\multicolumn{2}{@{}l}{\emph{Step~1} (\lean{Positivity.lean})} \\
$e^{Mt} \ge 0$ & \lean{exp_smul_apply_nonneg} \\
equation \eqref{eq:pattern} & \lean{exp_smul_apply_pos}, \lean{exp_smul_apply_eq_zero} \\
\end{tabular}\par
\vspace{0.4ex}
\noindent\begin{tabular}{@{}p{0.36\textwidth}p{0.60\textwidth}@{}}
\multicolumn{2}{@{}l}{\emph{Step~2} (\lean{Main.lean})} \\
equation \eqref{eq:support} & \lean{exp_smul_mulVec_pos}, \lean{exp_smul_mulVec_eq_zero} \\
the total is positive & \lean{sum_exp_smul_mulVec_pos} \\
equation \eqref{eq:noarc} & \lean{apply_eq_zero_of_mem_of_notMem} (\lean{Defs.lean}) \\
the reduced system & \lean{restrict_exp_smul_mulVec} \\
$B$ is Metzler & \lean{IsMetzler.reducedMatrix} (\lean{Defs.lean}) \\
\end{tabular}\par
\vspace{0.4ex}
\noindent\begin{tabular}{@{}p{0.36\textwidth}p{0.60\textwidth}@{}}
\multicolumn{2}{@{}l}{\emph{Step~3} (\lean{Modes.lean})} \\
equation \eqref{eq:expansion} & \lean{exp_smul_eq_sum_modes_card} \\
exponential decay beats any power of $t$ & \lean{tendsto_pow_mul_cexp_atTop_zero} \\
equation \eqref{eq:L} & \lean{exists_leading_term_real} \\
\end{tabular}\par
\vspace{0.4ex}
\noindent\begin{tabular}{@{}p{0.36\textwidth}p{0.60\textwidth}@{}}
\multicolumn{2}{@{}l}{\emph{Steps~4 to~6} (\lean{Main.lean})} \\
Step~4 (i) & \lean{leading_nonneg} \\
Step~4 (ii), equation \eqref{eq:semigroup} & \lean{exp_smul_mul_leading} \\
Step~4 (ii), equation \eqref{eq:eig} & \lean{mul_leading_eq} \\
Step~4 (iii) and Step~5, equation \eqref{eq:q} & \lean{tendsto_normalized_of_pos} \\
Step~6; Theorem~\ref{thm:main} & \lean{tendsto_normalized_solution} \\
\hline
\end{tabular}\par
}
\smallskip

\subsection{Where the formal proof differs from the text}\label{sec:lean:differences}

The formal proof follows the text step by step. It differs in the following points, none of which
changes the argument.

\begin{enumerate}[label=\textup{\arabic*.}, leftmargin=2em]
  \item \emph{Lemma~\ref{lem:exp} is not a single formal statement.} Each part is proved in the
        generality in which it is used. Parts~\ref{lem:exp:ode} and~\ref{lem:exp:semigroup} are
        proved for real matrices. Part~\ref{lem:exp:shift} is proved for a real matrix and a real
        $c$, and for a complex matrix in the form
        $e^{At} = e^{\lambda t} e^{(A - \lambda I)t}$. Part~\ref{lem:exp:eig} is proved for
        complex matrices. The statements about complex matrices allow $t$ to be complex.
  \item \emph{Real and complex matrices.} The text speaks of the eigenvalues of a real matrix
        $B$, meaning its complex eigenvalues. Lean keeps $B$ apart from the complex matrix with
        the same entries, written \verb|B.map (algebraMap ℝ ℂ)|. Two small lemmas connect them:
        \lean{exp_smul_map_ofReal} (the exponential of the second is that of the first, entry
        for entry) and \lean{ofReal_mem_spectrum_map_iff} (a real number is an eigenvalue of one
        exactly when it is an eigenvalue of the other).
  \item \emph{Lemma~\ref{lem:modes}.} In the formal statement the sum is over
        \verb|eigenvalues A|, the finite set of eigenvalues of $A$, and the multiplicity $m_\ell$
        is \verb|A.charpoly.rootMultiplicity lam|, the multiplicity of $\lambda_\ell$ as a root
        of the characteristic polynomial. The statement allows $t$ to be complex.
  \item \emph{Lemma~\ref{lem:PF}.} The formal statement is: there is a real number $\mu$ that is
        an eigenvalue of $B$ and such that every eigenvalue $\lambda \ne \mu$ of $B$ has
        $\Re\lambda < \mu$. That this $\mu$ is the largest real part of an eigenvalue is the
        separate statement \lean{isGreatest_re_spectrum}. The lemma assumes that $B$ has at least
        one row, which the text leaves tacit. In the proof of the theorem, $B = M_{RR}$ has at
        least one row because $S$, and hence $R$, is not empty.
  \item \emph{The constant $c$.} Where the text suggests $c = \max_i |M_{ii}|$, the formal proof
        takes $c = \sum_i |M_{ii}|$ (\lean{shiftConst}), in Step~1 and in the proof of
        Lemma~\ref{lem:PF}. All that matters is that $c \ge 0$ and $M + cI \ge 0$.
  \item \emph{Step~1.} The text describes $[P^k]_{ij}$ as a sum over sequences of indices. The
        formal proof does not expand the power. It proves the two implications by induction: if
        $[P^k]_{ij} \ne 0$ then $i$ is reachable from $j$, by induction on $k$
        (\lean{reachable_of_pow_apply_ne_zero}); and if $i$ is reachable from $j$ then
        $[P^k]_{ij} > 0$ for some $k$, by induction on the number of arcs
        (\lean{exists_pow_apply_pos_of_reachable}).
  \item \emph{Step~3.} The formal proof uses Lemma~\ref{lem:modes} with every inner sum running
        over $j < m$, the coefficients with $j \ge m_\ell$ being zero; equation
        \eqref{eq:expansion} leaves the range of $j$ unspecified. Also, the formal proof shows
        that $L$ is real already in this step, as a limit of real numbers, where the text notes
        it in Step~4~(i).
  \item \emph{Step~4~(ii).} Equation \eqref{eq:semigroup} is proved for every real $s$, not only
        for $s > 0$, and \eqref{eq:eig} is obtained by differentiating both sides at $s = 0$.
  \item \emph{Step~5.} The formal proof of the theorem takes $\tau = 1$.
\end{enumerate}

\subsection{What Mathlib supplies}\label{sec:lean:mathlib}

The formal proof takes the following from Mathlib: the matrix exponential, with its power series,
its derivative, and the rule $e^{X + Y} = e^{X} e^{Y}$ for commuting $X$ and $Y$ (in the formal
text, $e^{At}$ is written \verb|exp (t • A)|); the uniqueness theorem for solutions of
differential equations with a Lipschitz right-hand side, which is what justifies the sentence
beginning ``Hence'' in Lemma~\ref{lem:exp}\ref{lem:exp:ode}; facts (i) and (ii) about generalized
eigenspaces, used in the proof of Lemma~\ref{lem:modes}; the Cayley--Hamilton theorem; the
existence of an eigenvalue of a complex matrix with at least one row; and elementary analysis.
Fact (i) is Mathlib's \lean{Module.End.iSup_maxGenEigenspace_eq_top}, which its source labels
``Lemma~8.22(c)'' of \cite{Axler}, and fact (ii) is \lean{LinearMap.finrank_maxGenEigenspace_eq}.

Two things that this version of Mathlib does not have shaped the proofs in
Section~\ref{sec:background}.
\begin{itemize}[leftmargin=2em]
  \item \emph{The Perron--Frobenius theorem.} Mathlib defines irreducible and primitive matrices
        but proves no theorem about their eigenvalues. For this reason Lemma~\ref{lem:PF} is
        proved in full, from Step~1 and Lemma~\ref{lem:modes}, and not by reduction to
        \cite[Theorem~8.3.1]{HJ}.
  \item \emph{The Jordan normal form.} For this reason the proof of Lemma~\ref{lem:modes} uses
        generalized eigenspaces and not Jordan blocks.
\end{itemize}

\subsection{What has been checked, and what has not}\label{sec:lean:checked}

\begin{itemize}[leftmargin=2em]
  \item The project compiles without errors, and no file contains \lean{sorry}, Lean's
        placeholder for a missing proof.
  \item The main theorem and its three variants depend on three axioms only: \lean{propext},
        \lean{Classical.choice} and \lean{Quot.sound}. These are the standard axioms of
        classical mathematics in Lean, and Mathlib as a whole relies on them.
  \item The compiled files have been checked again with the program \lean{leanchecker}, which
        passes every definition and every proof through Lean's kernel once more.
\end{itemize}

These checks establish that the formal statement of Section~\ref{sec:lean:main} is a theorem.
They do not establish that the formal statement says what Theorem~\ref{thm:main} says. That is a
matter of reading the definitions and the statement against Section~\ref{sec:setting}, and
Sections~\ref{sec:lean:defs} and~\ref{sec:lean:main} are meant to make that reading possible.

\section*{Contributions}

The second author, Claude Fable~5.2 (Anthropic), is a large language model. Two separate
instances of the model took part in this work. They shared no memory, and each knew of the
other's work only through this text.

The first author, Thomas M.~English, posed the problem. The proof of Theorem~\ref{thm:main} was developed in a
dialogue between him and one instance of the model, which wrote the first draft of the text.

A second instance of the model wrote the Lean formalization described in Section~\ref{sec:lean}, under the
first author's direction. In the course of that work it rewrote the proofs of
Lemmas~\ref{lem:modes} and~\ref{lem:PF}, so that the text and the formal proof follow the same
argument: the proof of Lemma~\ref{lem:modes} now uses generalized eigenspaces where the original
used the Jordan normal form, and Lemma~\ref{lem:PF} is proved in full where the original cited
\cite[Theorem~8.3.1]{HJ}. It also wrote Section~\ref{sec:lean}.

The first author directed the work throughout, decided what was to be proved and how the result
was to be stated. He also reviewed the text, and made some minor revisions.
Responsibility for the content of this document, and for verification that the mechanically proven formal statement says what Theorem~\ref{thm:main} says (Section~\ref{sec:lean:checked}), lies entirely with him.
% <<<<<<<<<<<<<<<<<<<<<<<<<<<<<<<<<<<<<<<<<<<<<<<<<<<<<<<<<<<<<<<<<<<<<<<<<<<<<<<<<<<<<<<<<<<<<<<<<<<
% END OF LEAN NOTES
% <<<<<<<<<<<<<<<<<<<<<<<<<<<<<<<<<<<<<<<<<<<<<<<<<<<<<<<<<<<<<<<<<<<<<<<<<<<<<<<<<<<<<<<<<<<<<<<<<<<

% =====================================================================
\begin{thebibliography}{9}

\bibitem{Axler}
S.~Axler,
\emph{Linear Algebra Done Right}, 4th ed.,
Undergraduate Texts in Mathematics, Springer, 2024.

\bibitem{HJ}
R.~A. Horn and C.~R. Johnson,
\emph{Matrix Analysis}, 2nd ed.,
Cambridge University Press, Cambridge, 2013.

\end{thebibliography}

\end{document}
